\chapter{Estrategias de armonización}\label{cap:armonizacion}

\objetivos{
\begin{itemize}
  \item Elegir y justificar una unidad de conteo, una cuadrícula y un conjunto de satélites.
  \item Construir celdas-día con fuego a partir de una tabla de focos.
  \item Estimar razones de calibración por estratos y estabilizarlas cuando hay pocos datos.
  \item Encadenar calibraciones para cubrir todo el registro y propagar su incertidumbre.
  \item Diagnosticar si una serie armonizada es creíble y enunciar sus limitaciones.
\end{itemize}}

\section{Cuatro decisiones}

Todo lo anterior converge aquí. Armonizar los dos registros exige tomar cuatro decisiones, en
este orden, y cada una se apoya en capítulos anteriores (figura~\ref{fig:decisiones}).

\begin{figure}[ht]
\centering
\begin{tikzpicture}[node distance=6mm]
  \node[caja,fill=cieloClaro,text width=3.1cm] (a) {\textbf{1. Qué se cuenta}\\ unidad de actividad};
  \node[caja,fill=cieloClaro,text width=3.1cm,right=of a] (b) {\textbf{2. Dónde y cuándo}\\ cuadrícula y periodo};
  \node[caja,fill=cieloClaro,text width=3.1cm,right=of b] (c) {\textbf{3. Con qué ojos}\\ satélites y filtros};
  \node[caja,fill=bosqueClaro,text width=3.1cm,right=of c] (d) {\textbf{4. Cómo se enlaza}\\ calibración en el solape};
  \draw[flujo] (a)--(b); \draw[flujo] (b)--(c); \draw[flujo] (c)--(d);
\end{tikzpicture}
\caption{Las cuatro decisiones de una armonización.}
\label{fig:decisiones}
\end{figure}

Las tres primeras buscan que los dos sensores midan, en lo posible, \emph{la misma cosa}. Solo
la cuarta es estadística. Cuanto mejor se resuelvan las primeras, menos trabajo le queda a la
calibración y más creíble es el resultado.

\section{Primera decisión: la unidad de conteo}

Hay tres candidatas.

\begin{description}[leftmargin=1.2em,style=nextline]
  \item[Número de focos.] La más directa y la peor. Mezcla actividad con tamaño de píxel, número
  de satélites y duplicados (tabla~\ref{tab:multiple}).
  \item[Celdas-día con fuego.] Se superpone una cuadrícula fija y se cuenta, para cada celda y
  cada día, si hubo al menos una detección.
  \item[Resumen de potencia radiativa.] Puede resumirse la FRP por pasada comparable.
  Sumar potencias de instantes distintos no produce energía ni corrige el número de
  oportunidades de observación; hay que definir el muestreo temporal.
\end{description}

\begin{definicion}{celda-día con fuego}
Dada una cuadrícula fija, una celda-día con fuego es un par (celda, día local) en el que el
sensor registró al menos una detección. La actividad de una región en un periodo es el número de
celdas-día con fuego que contiene.
\end{definicion}

\begin{ejemplo}{de focos a celdas-día}
La figura~\ref{fig:celdadia} muestra las detecciones de un mismo día en una zona cubierta por
una cuadrícula de $4\times4$ celdas. Cuenta focos y celdas-día para cada sensor.
\solucion MODIS registró 4 focos, que caen en 3 celdas distintas. VIIRS registró 14 focos, en 5
celdas distintas.
\begin{center}
\begin{tabular}{lccc}
\toprule
 & MODIS & VIIRS & Razón VIIRS/MODIS \\
\midrule
Focos & 4 & 14 & 3,5 \\
Celdas-día con fuego & 3 & 5 & 1,7 \\
\bottomrule
\end{tabular}
\end{center}
Al cambiar de unidad, la razón entre sensores cae de 3,5 a 1,7. Lo que desapareció es la parte
de la diferencia que no era información: el incendio grande de la esquina, que VIIRS reparte en
ocho focos y MODIS en dos, cuenta ahora una celda para ambos. Lo que queda (dos celdas que solo
ve VIIRS) puede reflejar sensibilidad, cobertura o momento de observación; no se atribuye
toda la diferencia a una sola causa sin comprobarla.
\end{ejemplo}

\begin{atencion}[Cuánto reduce de verdad: el caso de Perú]
El ejemplo anterior está construido para que el efecto se vea. Con datos reales la reducción
puede ser mucho menor. En Perú, entre 2018 y 2020 y con los filtros del
capítulo~\ref{cap:calendario}, Suomi NPP registró 5,5 veces más focos diurnos que Aqua
(\num{201257} frente a \num{36546}); en celdas-día de $0{,}05^\circ$ la razón es
\PeruRatio{} (\PeruAjusteVIIRS{} frente a \PeruAjusteAqua{}). El cambio de unidad recorta
la diferencia solo un 15\,{}\%. En un país donde predominan las quemas pequeñas, casi toda
la ventaja de VIIRS consiste en fuegos que MODIS no registra, no en contar varias veces el
mismo. La celda-día sigue siendo la unidad correcta, pero no hace por sí sola el trabajo de
la calibración.
\end{atencion}

\begin{figure}[ht]
\centering
\begin{tikzpicture}[x=1.15cm,y=1.15cm]
  \foreach \sx/\tit/\col in {0/MODIS/modis,6.2/VIIRS/viirs}{
    \begin{scope}[shift={(\sx,0)}]
      \node[font=\small\bfseries,\col] at (2,4.4) {\tit};
      \ifdim\sx pt=0pt
        \foreach \i/\j in {0/3,2/1,3/0}{\fill[\col!25] (\i,\j) rectangle ++(1,1);}
        \foreach \x/\y in {0.4/3.5,0.75/3.3,2.5/1.45,3.4/0.5}{\fill[\col] (\x,\y) circle (0.09);}
      \else
        \foreach \i/\j in {0/3,2/1,3/0,1/1,3/2}{\fill[\col!25] (\i,\j) rectangle ++(1,1);}
        \foreach \x/\y in {0.2/3.7,0.35/3.55,0.5/3.4,0.65/3.25,0.8/3.1,0.3/3.3,0.55/3.65,0.75/3.5,
            2.45/1.5,2.6/1.35,3.4/0.5,3.55/0.4,1.5/1.6,3.6/2.4}{\fill[\col] (\x,\y) circle (0.07);}
      \fi
      \foreach \k in {0,...,4}{\draw[ceniza] (\k,0)--(\k,4); \draw[ceniza] (0,\k)--(4,\k);}
    \end{scope}
  }
\end{tikzpicture}
\caption{Las detecciones de un día en una misma zona (ejemplo ilustrativo). Los puntos son
focos; las celdas sombreadas son celdas-día con fuego.}
\label{fig:celdadia}
\end{figure}

La celda-día reduce la multiplicidad cuando varias detecciones comparten celda y día.
Puede reducir la sobredispersión, pero no elimina dependencias ni cambios de celda
por geolocalización. Añadir satélites sigue aumentando la probabilidad de registrar
alguna detección, aunque se tome la unión. La interpretación precisa es <<días con
alguna detección en esta celda bajo esta selección de plataformas>>.

Tiene también un coste: \textbf{satura}. Una celda con una quema y una celda con veinte cuentan
igual. Cuanto mayor es la celda, más actividad distinta queda fundida en un solo <<sí>>.

\section{Segunda decisión: la cuadrícula}

\subsection{Cuadrícula geográfica regular}

La cuadrícula más simple divide el mapa en celdas de lado fijo en grados de latitud y longitud.
Si el lado es $\Delta$ grados y el origen está en la latitud $-90^\circ$ y la longitud
$-180^\circ$, los índices de la celda que contiene un punto son
\begin{equation}\label{eq:indice}
  i = \left\lfloor \frac{\text{lat} + 90}{\Delta} \right\rfloor,
  \qquad
  j = \left\lfloor \frac{\text{lon} + 180}{\Delta} \right\rfloor,
\end{equation}
donde $\lfloor\cdot\rfloor$ es la parte entera hacia abajo. Es una operación trivial, que se
aplica a millones de filas en una fracción de segundo.

¿Cuánto mide una celda en kilómetros? Un grado de latitud equivale siempre a unos 111,3~km. Un
grado de longitud equivale a esa misma distancia en el ecuador y se estrecha hacia los polos:
\begin{equation}\label{eq:grado}
  1^\circ\text{ de latitud} \approx \SI{111.3}{\kilo\metre},
  \qquad
  1^\circ\text{ de longitud} \approx \SI{111.3}{\kilo\metre}\times\cos(\text{lat}) .
\end{equation}

\begin{ejemplo}{una celda sobre Huancayo}
Para una cuadrícula de $\Delta = 0{,}05^\circ$, halla los índices de la celda que contiene a
Huancayo (latitud $-12{,}07^\circ$, longitud $-75{,}21^\circ$) y sus dimensiones.
\solucion Índices: $i = \lfloor (-12{,}07 + 90)/0{,}05 \rfloor = \lfloor 1558{,}6 \rfloor =
1558$; $j = \lfloor (-75{,}21 + 180)/0{,}05 \rfloor = \lfloor 2095{,}8 \rfloor = 2095$.
Dimensiones: de norte a sur, $0{,}05 \times 111{,}3 = \SI{5.57}{\kilo\metre}$; de este a oeste,
$5{,}57 \times \cos(12{,}07^\circ) = 5{,}57 \times 0{,}978 = \SI{5.44}{\kilo\metre}$. La celda
mide unos $5{,}6 \times 5{,}4$~km, casi \SI{30}{\kilo\metre\squared}.
\end{ejemplo}

En las latitudes del Perú la deformación es pequeña (el coseno va de 1,00 en el ecuador a 0,95
en Tacna), por lo que una rejilla en grados puede ser una aproximación práctica. Para
comparar densidades se consideran las áreas reales. Para un producto global
habría que compensar el estrechamiento de las celdas hacia los polos o usar una cuadrícula de
igual área.

\subsection{El tamaño de la celda}

Es la decisión con más consecuencias y no tiene una respuesta única. La
figura~\ref{fig:tamano-celda} resume el compromiso.

\begin{figure}[ht]
\centering
\begin{tikzpicture}
\begin{axis}[width=13cm,height=6.6cm,xmode=log,xmin=0.3,xmax=120,ymin=0.8,ymax=4.2,
  xlabel={Lado de la celda (km)},ylabel={Razón VIIRS/MODIS en celdas-día},
  xtick={0.375,1,5,10,25,50,100},xticklabels={{0,4},1,5,10,25,50,100},
  grid=major,grid style={ceniza!15},samples=100,domain=0.3:120]
\addplot[fuego,very thick] {1+2.6/(1+(x/4)^0.9)};
\addplot[ceniza,dashed] {1};
\fill[bosque,opacity=0.13] (axis cs:4,0.8) rectangle (axis cs:12,4.2);
\node[font=\footnotesize,align=center] at (axis cs:7,3.75) {intervalo\\ razonable};
\node[font=\footnotesize,align=left,anchor=west,fill=white,inner sep=2pt] at (axis cs:0.36,1.7) {celda más fina que\\ el píxel de MODIS:\\ razón cercana a\\ la de los focos};
\node[font=\footnotesize,align=right,anchor=east] at (axis cs:110,2.4) {celda enorme:\\ casi todo satura,\\ se pierde detalle};
\end{axis}
\end{tikzpicture}
\caption{Cómo cambia la razón entre sensores con el tamaño de la celda (curva esquemática). Con
celdas muy pequeñas la celda-día equivale a contar focos; con celdas muy grandes ambos sensores
dicen <<sí>> casi siempre y la razón tiende a 1, a costa de borrar toda la estructura espacial.}
\label{fig:tamano-celda}
\end{figure}

Tres criterios acotan la elección:

\begin{itemize}
  \item \textbf{Relacionada con la huella.} Conviene estudiar escalas mayores que la
  dimensión gruesa de MODIS lejos del nadir (casi 5~km). Aun así, un fuego junto al borde
  puede cambiar de celda; ninguna rejilla elimina ese efecto por completo.
  \item \textbf{Suficientes datos por celda.} Por la ecuación~\eqref{eq:rel}, una celda con
  pocas celdas-día al año tiene estimaciones muy inciertas.
  \item \textbf{Adecuada al uso.} Un gestor provincial no necesita más detalle que unos
  kilómetros.
\end{itemize}

Un lado entre 0,05$^\circ$ y 0,1$^\circ$ (aproximadamente 5--11~km) es un punto de partida
para comparar variantes regionales, no una resolución validada de antemano. Se estudia
cómo cambian los indicadores al variar el tamaño y se documenta su unidad.

\begin{parareto}
Guarda los datos en la cuadrícula \emph{fina} (por ejemplo 0,05$^\circ$) y agrega hacia arriba
cuando el usuario elija un área: distrito, provincia, región. La calibración, en cambio, se
estima en unidades grandes, donde hay datos suficientes. Son dos resoluciones distintas para dos
tareas distintas, y confundirlas es un error frecuente.
\end{parareto}

\subsection{La resolución temporal}

La celda-día fija el día como unidad mínima; aquí elegimos el \emph{día local}
(capítulo~\ref{cap:firms}). Para presentar, se agrega a semanas o a meses. Hay una sutileza: las
semanas no encajan en los meses ni en los años. Una opción cómoda es dividir cada año en
periodos fijos de ocho días (46 por año) o trabajar con el \emph{día del año} (de 1 a 365) y
aplicar una media móvil.

\section{Tercera decisión: satélites y filtros}

Con lo visto en los capítulos~\ref{cap:orbitas} y~\ref{cap:errores}, la configuración que
produce los dos registros más parecidos es:

\begin{center}
\renewcommand{\arraystretch}{1.2}
\begin{tabularx}{\textwidth}{l>{\raggedright\arraybackslash}X>{\raggedright\arraybackslash}X}
\toprule
 & Era MODIS & Era VIIRS \\
\midrule
Satélite & Aqua & Suomi NPP \\
Motivo & \multicolumn{2}{>{\raggedright\arraybackslash}p{10cm}}{horas nominales próximas (13:30 y 01:30); verificar pasadas reales y deriva} \\
Tipo & \texttt{type = 0} & \texttt{type = 0} \\
Confianza & $\geq 30$ & nominal y alta \\
Día y noche & por separado & por separado \\
Colección & 6.1 & 2 \\
\bottomrule
\end{tabularx}
\end{center}

Queda un problema: Aqua no tiene datos antes de mediados de 2002. Para los primeros años solo
existe Terra, que observa a otra hora. La solución es otra calibración, entre Terra y Aqua, en
\emph{su} periodo común, que es larguísimo. La serie completa se arma entonces en tres tramos
enlazados (figura~\ref{fig:cadena}).

\begin{figure}[ht]
\centering
\begin{tikzpicture}[x=0.5cm,y=0.75cm]
  \draw[->,thick] (0,0)--(27.5,0);
  \foreach \a in {2000,2004,2008,2012,2016,2020,2024}{\pgfmathsetmacro{\x}{\a-2000}\draw (\x,0.12)--(\x,-0.12) node[below,font=\footnotesize]{\a};}
  \fill[modis!85] (0.85,3.6) rectangle (26.5,4.1); \node[left,font=\small] at (0.6,3.85) {Terra};
  \fill[modis!55] (2.5,2.3) rectangle (26.5,2.8); \node[left,font=\small] at (0.6,2.55) {Aqua};
  \fill[viirs!80] (12.05,1.0) rectangle (26.5,1.5); \node[left,font=\small] at (0.6,1.25) {Suomi NPP};
  \draw[thick,<->,>={Stealth[length=1.3mm]}] (5,2.9)--(5,3.5)
    node[midway,left,font=\footnotesize] {$r_1$};
  \draw[thick,<->,>={Stealth[length=1.3mm]}] (15,1.6)--(15,2.2)
    node[midway,right,font=\footnotesize] {$r_2$};
  \node[font=\footnotesize,align=center,text width=5.4cm] at (6.5,5.1)
    {\textbf{Calibración $r_1$}\\Terra $\to$ Aqua, 2003--2011};
  \node[font=\footnotesize,align=center,text width=5.4cm] at (20,5.1)
    {\textbf{Calibración $r_2$}\\Aqua $\to$ Suomi NPP, 2012--2020};
  \draw[fuego,line width=2.2pt] (0.85,-1.6)--(2.5,-1.6); \draw[brasa,line width=2.2pt] (2.5,-1.6)--(12.05,-1.6); \draw[bosque,line width=2.2pt] (12.05,-1.6)--(26.5,-1.6);
  \node[font=\footnotesize,fuego,align=center] at (1.6,-2.4) {Terra\\ $\times r_1 r_2$};
  \node[font=\footnotesize,brasa!70!black] at (7.2,-2.2) {Aqua $\times\, r_2$};
  \node[font=\footnotesize,bosque] at (19,-2.2) {Suomi NPP, tal cual};
  \node[left,font=\small\bfseries] at (0.6,-1.6) {Serie};
\end{tikzpicture}
\caption{Calibración encadenada. Para expresar toda la serie en la escala de Suomi NPP, los años
de Aqua se multiplican por $r_2$ y los años en que solo hay Terra por $r_1 r_2$.
Los periodos ilustran una estrategia; la calibración real peruana del
capítulo~\ref{cap:aplicacion} utiliza un intervalo de ajuste distinto.}
\label{fig:cadena}
\end{figure}

Cada eslabón añade incertidumbre. Si las incertidumbres relativas de $r_1$ y de $r_2$ son
$u_1$ y $u_2$ y son independientes, la del producto $r_1 r_2$ es
\begin{equation}\label{eq:cuadratura}
  u_{12} \simeq \sqrt{u_1^2 + u_2^2} .
\end{equation}

\begin{ejemplo}{incertidumbre de una cadena}
Las razones se estimaron con incertidumbres estándar relativas del 5\,{}\% para $r_1$ y del 4\,{}\% para
$r_2$. ¿Cuál es la incertidumbre de calibración de los años 2000--2002?
\solucion $u_{12} = \sqrt{0{,}05^2 + 0{,}04^2} = \sqrt{0{,}0041} = 0{,}064$, es decir, un
6,4\,{}\%. Los años más antiguos son los de mayor incertidumbre de calibración, y además los de
menos observaciones. Conviene mostrarlos con una banda de error más ancha, o incluso empezar la
serie en 2003 y decirlo.
\end{ejemplo}

\begin{atencion}[Terra y Aqua no difieren solo en un factor]
La pasada de las 10:30 ve otro momento del ciclo diurno, y la proporción entre lo que arde a
media mañana y a primera hora de la tarde depende del tipo de fuego. La razón $r_1$ varía con la
región y la estación tanto o más que $r_2$. Si no hay tiempo de tratarla con cuidado, es
legítimo y más honesto limitar la serie a 2003 en adelante.
\end{atencion}

\section{Cuarta decisión: la calibración en el solape}

\subsection{En qué dirección convertir}

Hay que elegir una escala común, y hay dos opciones.

\begin{description}[leftmargin=1.2em,style=nextline]
  \item[Llevar MODIS a la escala de VIIRS.] Se multiplican los conteos antiguos por la razón.
  La serie queda en una escala VIIRS especificada. El cambio posterior de plataforma
  necesita otra evaluación: Suomi NPP dejó de entregar datos nuevos en noviembre de 2026.
  Los años antiguos se convierten en una estimación de actividad equivalente que
  MODIS nunca observó: se supone que los fuegos pequeños guardaban entonces la misma proporción
  con los grandes que durante el solape.
  \item[Llevar VIIRS a la escala de MODIS.] Se transforman los conteos recientes en un
  índice equivalente, con incertidumbre y posible pérdida de sensibilidad, ligado a la
  escala MODIS. Tampoco es una nueva observación directa.
\end{description}

Ambas direcciones requieren validación. Una constante común aplicada a un valor y a su
propia referencia se cancela en una anomalía relativa; esto no garantiza equivalencia
si la referencia mezcla épocas o las razones cambian por estrato. El objetivo no es
forzar continuidad, sino reducir discrepancias de observación sin borrar cambios reales.
Aquí se ilustra la conversión a una escala VIIRS, declarando sus supuestos.

\subsection{Estratificar}

La lección repetida de los capítulos~\ref{cap:modisviirs} y~\ref{cap:deteccion} es que la razón
entre sensores depende del tipo de fuego. La forma práctica de recogerlo es
\termino{estratificar}: dividir los datos en grupos donde el tipo de fuego sea homogéneo y
estimar una razón para cada uno.

\begin{definicion}{estrato}
Un estrato es un subconjunto de los datos definido por variables que se sospecha que influyen en
la relación que se quiere estimar. Aquí, combinaciones de zona geográfica, época del año y día o
noche.
\end{definicion}

Para el Perú, una estratificación geográfica natural son las tres grandes regiones, costa,
sierra y selva, o los departamentos. En el tiempo, los meses o los trimestres. La
tabla~\ref{tab:estratos} muestra el aspecto que tiene el resultado.

\begin{table}[ht]
\centering
\begin{tabular}{lrrrr}
\toprule
 & \multicolumn{4}{c}{Razón Suomi NPP / Aqua, celdas-día diurnas} \\
\cmidrule{2-5}
Zona & ene--mar & abr--jun & jul--sep & oct--dic \\
\midrule
Sierra central & (2,9) & 2,4 & 1,9 & 2,1 \\
Selva central & (2,2) & 1,9 & 1,5 & 1,6 \\
Costa & (3,4) & (3,1) & 2,8 & (3,0) \\
\bottomrule
\end{tabular}
\caption{Aspecto de una tabla de calibración estratificada (valores ilustrativos, no medidos).
Entre paréntesis, estratos con tan pocos datos que su razón no es fiable.}
\label{tab:estratos}
\end{table}

Una tabla así, con valores reales, ya es un resultado interesante: la razón baja en plena
temporada, cuando hay fuegos mayores, y sube fuera de ella. Es la física del
capítulo~\ref{cap:deteccion} hecha visible.

Los valores de la tabla son inventados y, para Perú, demasiado bajos. Con datos reales
(sección~\ref{sec:peru-modelos}) la razón nacional mensual entre celdas-día de Suomi NPP y de
Aqua va de \PeruRazonMin{} en septiembre a \PeruRazonMax{} en marzo. El patrón, en cambio, es
el previsto: mínima en plena temporada y máxima fuera de ella.

\subsection{Estratos con pocos datos}

Estratificar tiene un precio. Cuantos más estratos, menos datos en cada uno, y una razón
calculada con seis celdas-día es ruido (ecuación~\eqref{eq:rel}). Hay dos remedios, que pueden
combinarse.

\textbf{Retroceso jerárquico.} Se fija un mínimo de datos, por ejemplo 100 celdas-día de MODIS
en el estrato. Si un estrato no lo alcanza, se usa la razón del nivel superior: del trimestre al
año, del departamento a la región natural, de la región al país.

\textbf{Contracción hacia la razón general.} En lugar de un salto brusco, se mezcla
gradualmente la razón del estrato con la general:
\begin{equation}\label{eq:contraccion}
  r_s = \frac{Y_s + k\,r_g}{X_s + k},
\end{equation}
donde $X_s$ e $Y_s$ son las sumas de MODIS y de VIIRS en el estrato, $r_g$ es la razón general y
$k$ es una constante que se elige (del orden de unas decenas). La fórmula equivale a añadir al
estrato un pseudoconteo $k$ en la escala MODIS, no necesariamente $k$ observaciones
independientes. Si el estrato
tiene muchos datos ($X_s \gg k$), manda su propia razón; si tiene pocos, manda la general.

\begin{ejemplo}{contracción de una razón}
Un estrato tiene $X_s = 6$ celdas-día de MODIS y $Y_s = 30$ de VIIRS. La razón general de su
región es $r_g = 2{,}9$. Calcula la razón cruda y la contraída con $k = 20$. Repite para un
estrato con $X_s = 600$ e $Y_s = 3000$.
\solucion Primer estrato: razón cruda $30/6 = 5{,}0$. Contraída: $(30 + 20\times2{,}9)/(6+20) =
88/26 = 3{,}4$. Con solo seis observaciones no se le da crédito al valor extremo. Segundo
estrato: razón cruda $5{,}0$; contraída $(3000 + 58)/(600+20) = 4{,}9$. Con seiscientas
observaciones, la razón propia se impone.
\end{ejemplo}

\subsection{Acercar los sensores antes de calibrar}

Una razón estadística compensa las diferencias en promedio. Otra vía, complementaria, es reducir
las diferencias \emph{físicamente}, descartando de VIIRS lo que MODIS no podría haber visto. La
celda-día ya lo hace en parte. Una segunda herramienta es un umbral de potencia.

Como el área mínima detectable es proporcional al área del píxel (ecuación~\eqref{eq:amin}), y
la FRP es proporcional al área y a la cuarta potencia de la temperatura, MODIS tiene un límite
inferior aproximado de FRP por debajo del cual casi nunca detecta. Si de VIIRS se eliminan las
celdas cuya FRP de una pasada comparable queda por debajo de un umbral $F_0$, podría
reducirse parte de la diferencia de sensibilidad. No se suman pasadas de todo un día
como si fueran simultáneas. El umbral se estudia con datos de ajuste y se evalúa fuera
de muestra; una razón cercana a 1 por sí sola no demuestra acuerdo entre detecciones
(figura~\ref{fig:umbral}).

\begin{figure}[ht]
\centering
\begin{tikzpicture}
\begin{axis}[width=12.5cm,height=6.2cm,xmin=0,xmax=30,ymin=0.6,ymax=2.6,
  xlabel={Umbral de FRP aplicado a las celdas-día de VIIRS, $F_0$ (MW)},
  ylabel={Razón VIIRS filtrado / MODIS},
  legend style={at={(0.97,0.95)},anchor=north east,font=\small,draw=ceniza!40},
  grid=major,grid style={ceniza!15},samples=80,domain=0:30]
\addplot[bosque,very thick] {0.72+1.55*exp(-x/6.5)};
\addplot[brasa,very thick,dashed] {0.80+1.05*exp(-x/5)};
\addplot[ceniza,dashed] {1};
\legend{fuera de temporada,en temporada}
\end{axis}
\end{tikzpicture}
\caption{Efecto esperado de un umbral de potencia (curvas esquemáticas). Al subir el umbral, las
detecciones de VIIRS que MODIS no habría captado desaparecen y la razón se acerca a 1. Si se
sube demasiado, se eliminan también fuegos que MODIS sí ve.}
\label{fig:umbral}
\end{figure}

\begin{atencion}[Un método con supuestos]
El umbral de FRP no es un procedimiento estándar: es una idea razonable, basada en la física del
libro, cuya eficacia hay que comprobar con los datos. Sus supuestos son discutibles (la FRP de
VIIRS proviene de píxeles de 750~m; el límite de MODIS depende del ángulo de barrido). Si lo
usas, muéstralo como una variante comparada con la razón simple, no como una verdad.
\end{atencion}

\section{La receta completa}

La figura~\ref{fig:receta} reúne el procedimiento.

\begin{figure}[ht]
\centering
\begin{tikzpicture}[node distance=4mm,
  paso/.style={caja,text width=12cm,align=left,fill=cieloClaro}]
  \node[paso] (a) {\textbf{1. Limpiar} cada tabla de focos (figura~\ref{fig:limpieza}) y convertir a día local.};
  \node[paso,below=of a] (b) {\textbf{2. Asignar celda} a cada foco con la ecuación~\eqref{eq:indice}.};
  \node[paso,below=of b] (c) {\textbf{3. Reducir a presencias}: una fila por (plataforma, celda, día, franja); conservar por separado información de pasadas y cobertura.};
  \node[paso,below=of c] (d) {\textbf{4. Separar el solape} en años de ajuste y años de prueba.};
  \node[paso,below=of d] (e) {\textbf{5. Estimar razones} por estrato con los años de ajuste (razón de sumas, con contracción).};
  \node[paso,below=of e] (f) {\textbf{6. Validar}: comparar años de prueba y medir sesgo y MAE; usar métricas relativas con denominadores válidos.};
  \node[paso,below=of f] (g) {\textbf{7. Reestimar} las razones con todo el solape y construir la serie completa.};
  \node[paso,below=of g,fill=bosqueClaro] (h) {\textbf{8. Diagnosticar} la serie y adjuntar su incertidumbre.};
  \foreach \u/\v in {a/b,b/c,c/d,d/e,e/f,f/g,g/h}{\draw[flujo] (\u)--(\v);}
\end{tikzpicture}
\caption{Procedimiento de armonización de principio a fin.}
\label{fig:receta}
\end{figure}

\section{Diagnósticos: ¿es creíble la serie?}

Una serie armonizada debe superar tres comprobaciones antes de mostrarse a nadie.

\textbf{Discrepancias en el empalme.} Una ruptura coincidente con el cambio de sensor
merece investigación, pero puede coexistir con un cambio real de actividad. Se comparan
en el solape MODIS convertido y VIIRS observado, con sus errores y cobertura. No se
suaviza la serie solo para obligar a desaparecer un salto.

\textbf{Estabilidad de la razón en el tiempo.} Se calcula la razón VIIRS/MODIS de cada año del
solape. Si la calibración es válida, esos valores oscilan alrededor de una constante. Una
tendencia sostenida (figura~\ref{fig:estabilidad}) puede reflejar cambios de mezcla de
fuegos, cobertura, procesamiento u hora de paso; hay que investigar antes de extrapolar.

\begin{figure}[ht]
\centering
\begin{tikzpicture}
\begin{axis}[width=13cm,height=6cm,xmin=2011.5,xmax=2024.5,ymin=1.2,ymax=3.2,
  xlabel={Año},ylabel={Razón anual VIIRS/MODIS},
  xtick={2012,2014,2016,2018,2020,2022,2024},x tick label style={/pgf/number format/1000 sep=},
  legend style={at={(0.03,0.96)},anchor=north west,font=\small,draw=ceniza!40},
  grid=major,grid style={ceniza!15}]
\addplot[bosque,very thick,mark=*] coordinates {(2012,1.92)(2013,1.85)(2014,1.97)(2015,1.88)(2016,1.94)(2017,1.86)(2018,1.93)(2019,1.9)(2020,1.96)};
\addplot[fuego,very thick,mark=square*] coordinates {(2020,1.96)(2021,2.08)(2022,2.24)(2023,2.47)(2024,2.71)};
\legend{estable: calibración válida,tendencia: algo cambia en un sensor}
\draw[ceniza,dashed] (axis cs:2011.5,1.91)--(axis cs:2024.5,1.91);
\end{axis}
\end{tikzpicture}
\caption{Diagnóstico de estabilidad (datos ilustrativos). Mientras la razón anual fluctúa
alrededor de un valor fijo, este diagnóstico no detecta una deriva. Si cambia, se
investigan causas: la deriva orbital es una posibilidad y no fija por sí sola el signo. En
Perú, la razón anual real entre 2013 y 2024 oscila entre 4,0 y 5,5 sin tendencia clara
(\texttt{datos/peru\_anual.csv}): la variación de un año a otro es grande y está ligada a
cuánta actividad hubo.}
\label{fig:estabilidad}
\end{figure}

\textbf{Residuos sin estructura.} Los errores de la validación, agrupados por mes del año y por
zona, no deben mostrar un patrón. Si el método sobrestima siempre agosto y subestima siempre
abril, falta estratificar por estación.

\section{Limitaciones que hay que declarar}

Ninguna armonización es perfecta, y un trabajo serio enumera lo que no resuelve. Para la receta
de este capítulo:

\begin{enumerate}
  \item \textbf{Supone proporciones constantes.} Se asume que la mezcla de fuegos pequeños y
  grandes en la era MODIS era la misma que en el solape. Si las prácticas agrícolas cambiaron,
  la reconstrucción de los años antiguos es incorrecta, y los datos no permiten saberlo.
  \item \textbf{No recupera información.} Los fuegos que MODIS no vio se compensan en promedio,
  no se recuperan. La estimación solo es defendible en escalas donde se haya validado;
  no localiza retrospectivamente cada fuego omitido.
  \item \textbf{No corrige las nubes.} Un mes muy nuboso aparece como un mes de poca actividad.
  \item \textbf{Depende de las pasadas elegidas.} Las horas nominales pueden derivar y
  las pasadas no capturan toda la evolución del fuego.
  \item \textbf{La unidad satura.} Las celdas-día no distinguen una quema de muchas en la misma
  celda.
\end{enumerate}

\begin{parareto}
Estas cinco frases, puestas en una diapositiva o en la documentación, no debilitan el proyecto:
lo fortalecen. Muestran que el equipo entiende sus datos, que es exactamente lo que busca un
reto de armonización propuesto por quienes producen esos datos.
\end{parareto}

\resumen

\begin{itemize}
  \item Armonizar exige cuatro decisiones: \textbf{unidad}, \textbf{cuadrícula},
  \textbf{satélites y filtros}, y \textbf{calibración}.
  \item La \textbf{celda-día con detección} reduce multiplicidad dentro de una clave;
  no iguala sensibilidad ni cobertura y puede saturar.
  \item Una cuadrícula en grados se indexa con una división entera. Un lado de 0,05$^\circ$ a
  0,1$^\circ$ es razonable para uso regional.
  \item La pareja \textbf{Aqua / Suomi NPP} comparte hora de paso. Los años anteriores a 2002
  requieren una segunda calibración con Terra, y las incertidumbres se suman en cuadratura.
  \item La razón se estima \textbf{por estratos} de zona y época, con retroceso jerárquico o
  contracción cuando hay pocos datos.
  \item Un \textbf{umbral de FRP} puede acercar físicamente VIIRS a MODIS, con supuestos que hay
  que comprobar.
  \item Toda serie armonizada debe pasar tres diagnósticos: continuidad en el empalme,
  estabilidad de la razón y residuos sin estructura.
  \item Las limitaciones se declaran.
\end{itemize}

\ejercicios

\begin{enumerate}
  \item En una celda, un día, Suomi NPP registra 9 focos, NOAA-20 registra 11 y Aqua registra 2.
  ¿Cuántos focos y cuántas celdas-día aporta cada sensor? ¿Y si se juntan los dos VIIRS?
  \item Calcula los índices de celda, con $\Delta = 0{,}1^\circ$, de un foco en latitud
  $-6{,}48^\circ$ y longitud $-76{,}37^\circ$ (Tarapoto). ¿Cuánto mide esa celda en kilómetros?
  \item ¿Cuántas celdas de $0{,}05^\circ$ cubren aproximadamente el departamento de Junín (unos
  \SI{44000}{\kilo\metre\squared})? Si en un año hay 3000 celdas-día con fuego en el
  departamento, ¿cuántas tocan en promedio a cada celda? ¿Qué te dice eso sobre calcular
  estadísticas celda a celda?
  \item Explica con la figura~\ref{fig:tamano-celda} por qué no conviene una celda de 500~m ni
  una de 100~km.
  \item Una cadena de calibración tiene tres eslabones con incertidumbres relativas del 3\,{}\%,
  el 5\,{}\% y el 6\,{}\%. Calcula la incertidumbre del producto.
  \item Aplica la ecuación~\eqref{eq:contraccion} con $k = 30$ y $r_g = 2{,}0$ a tres estratos:
  $(X, Y) = (3, 15)$, $(40, 100)$ y $(900, 1500)$. Comenta.
  \item Para un calendario que solo muestra anomalías relativas dentro de cada época, ¿importa
  si se convierte MODIS a escala VIIRS o al revés? ¿Y para una gráfica de la tendencia de 2003 a
  hoy?
  \item La razón anual VIIRS/MODIS de tu región es estable de 2012 a 2019 y crece un 8\,{}\% cada
  año desde 2021. Propón una causa, una forma de comprobarla con las columnas de los datos y una
  decisión sobre qué años usar para calibrar.
\end{enumerate}

\lecturas

Los artículos de los algoritmos y las guías explican las diferencias físicas entre
sensores.\fuente{giglio2016}\fuente{schroeder2014} La receta de razones y celdas-día
es una propuesta didáctica que debe validarse, no un procedimiento oficial tomado
de esos artículos. Los métodos de continuidad de área quemada no se trasladan sin
más a fuego activo. El capítulo~\ref{cap:aplicacion} prueba una línea base real en Perú.
